\chapter{Conclusão}
\label{chap:conclusao}

Este trabalho analisou a reação fiscal dos estados brasileiros perante o endividamento e se tal resposta é condicionada pelo desenho dos contratos de refinanciamento da Lei n\textsuperscript{o}~9.496/1997. Adaptando a função de reação fiscal de \citeonline{abubakarDebtReliefFiscal2025} ao contexto subnacional, estimou-se por 2SLS uma especificação que interage a dívida defasada com o teto contratual de comprometimento da receita corrente líquida (RCL).

Os resultados sustentam duas conclusões. Primeiro, a reação fiscal líquida --- combinação dos coeficientes da dívida defasada e de sua interação com o teto --- é positiva e estatisticamente significativa em todos os grupos de teto observados, indicando que os estados elevam o resultado primário quando o endividamento cresce e satisfazendo a condição de sustentabilidade de \citeonline{bohn1998}. Segundo, o coeficiente negativo da interação mostra que a intensidade dessa reação está associada ao desenho contratual: estados com tetos mais permissivos ajustam menos. A reação fiscal líquida varia de aproximadamente 0,15 nos estados com teto de 12\% a 0,03 naqueles com teto de 15\%, permanecendo positiva em todos os grupos, ainda que apenas marginalmente significativa no grupo de teto mais elevado. O exercício de robustez com a tripla interação de $binding$ não fornece evidência robusta de heterogeneidade associada à proximidade do teto: o coeficiente é apenas marginalmente significativo, de sinal contrário ao esperado e identificado a partir de pouquíssimas observações. O efeito moderador opera, assim, pelo nível contratual em si, e não pelo risco iminente de ultrapassá-lo.

A interpretação dos resultados, no entanto, exige cautela, uma vez que renegociações posteriores à lei original --- a troca do indexador pela LC n\textsuperscript{o}~148/2014 e a extensão de prazo pela LC n\textsuperscript{o}~156/2016, por exemplo --- reduziram o serviço efetivo da dívida e afastaram a maioria dos estados do teto contratual, confirmado pela baixa frequência do \textit{binding} na amostra. O coeficiente $\hat{\beta}_2$ deveria, portanto, ser lido menos como efeito de uma restrição mecânica ativa e mais como uma medida da heterogeneidade institucional no desenho original dos contratos, associada a diferenças no comportamento fiscal ao longo de todo o período amostral. A robustez à exclusão dos estados sob suspensão de pagamentos (Rio de Janeiro, Rio Grande do Sul e Goiás) reforça essa leitura, ao mostrar que o resultado não decorre desses episódios de estresse fiscal agudo.

Esse achado refina o resultado de \citeonline{abubakarDebtReliefFiscal2025}: enquanto na África Subsaariana a presença de regras fiscais pode ser contraproducente, no Brasil regras mais permissivas atenuam o ajuste sem invertê-lo. A divergência, conforme discutido no capítulo anterior, é atribuível ao arcabouço institucional brasileiro, em particular ao papel da União como credora direta e ao \textit{enforcement} doméstico via Lei de Responsabilidade Fiscal (LRF). A distinção relevante para o desenho de políticas, portanto, não é a mera existência de regras fiscais, mas a calibração dos parâmetros contratuais e o arcabouço institucional que as sustenta. Investigações futuras poderiam avaliar se o efeito aqui identificado se estende a outras dimensões da política fiscal, como a composição do gasto público, ou como se manifesta nas rodadas de renegociação em curso.

